\chapter{Strutture dati e tipi di dato astratti}
\label{chap:strutture-dati-adt}

Gli algoritmi di ricerca e ordinamento del \cref{chap:ricerca-ordinamento}
operano su collezioni di dati. La scelta di come organizzarle dipende
sia dalle operazioni richieste dal problema, sia dal costo della loro
esecuzione, studiato nel \cref{chap:complessita}. Si distinguono pertanto
il comportamento del contenitore e la sua rappresentazione interna.\footnote{Materiale
di riferimento: G. Baruzzo, \emph{Strutture dati e tipi di dati astratti},
dispense del corso di Fondamenti di Informatica.}

\section{Tipi di dato astratti e strutture dati}
\label{sec:fi-adt-strutture}

\begin{definition}[Tipo di dato astratto]
\label{def:fi-adt}
Un \emph{tipo di dato astratto} (\emph{Abstract Data Type}, ADT)
specifica le operazioni su una collezione di dati e il loro
comportamento, compresi risultati, precondizioni e casi di errore,
senza imporne l'implementazione.
\end{definition}

Una \emph{struttura dati} stabilisce invece come gli elementi e le
informazioni necessarie a collegarli siano organizzati in memoria.
Più strutture possono realizzare lo stesso ADT, con costi temporali
e spaziali differenti. Per esempio, una coda di richieste deve
elaborare gli arrivi nell'ordine in cui si presentano: questo vincolo
definisce il comportamento richiesto, ma non stabilisce come
memorizzare le richieste.

Nella progettazione si individua prima l'ADT adatto al problema,
quindi se ne sceglie un'implementazione in base alle operazioni più
frequenti.

\section{Array e accesso diretto}
\label{sec:fi-strutture-array}

Un \emph{array} memorizza \(n\) elementi di dimensione fissa
\(s\) in posizioni contigue. Indicando con \(b\) l'indirizzo del primo
elemento, l'indirizzo dell'elemento di indice \(i\), per
\(0\le i<n\), è
\[
  \operatorname{addr}(v[i])=b+i\,s.
\]
Il calcolo richiede un numero costante di operazioni: accesso e
aggiornamento per indice hanno costo \(\Theta(1)\), nel modello della
\cref{sec:fi-complessita-modello}. Gli elementi occupano \(ns\)
byte, ai quali si aggiungono le informazioni di gestione.

L'organizzazione contigua favorisce la località spaziale descritta
nella \cref{sec:fi-memoria-archiviazione}. La \emph{lista} di Python
della \cref{sec:fi-python-liste} è realizzata, in CPython, mediante
un array dinamico di riferimenti agli oggetti: sono contigui i
riferimenti, non necessariamente gli oggetti stessi.

\subsection{Array statici e dinamici}
\label{sec:fi-array-dinamici}

Un array \emph{statico} ha lunghezza fissata; un array \emph{dinamico}
può variare il numero degli elementi. Per quest'ultimo si distinguono
la \emph{dimensione} \(n\), cioè il numero di elementi presenti, e la
\emph{capacità} \(C\ge n\), cioè i posti disponibili nell'area
allocata. Quando \(n=C\), un inserimento può richiedere un'area più
grande, la copia degli \(n\) elementi e il rilascio della precedente:
il singolo ridimensionamento costa \(\Theta(n)\).

Se la capacità viene aumentata di un fattore costante maggiore di
uno, per esempio raddoppiata, il costo complessivo delle copie
su \(m\) inserimenti in coda è limitato da una somma geometrica:
\[
  1+2+4+\cdots+2^{\lfloor\log_2 m\rfloor}<2m.
\]
Il costo totale di \(m\) inserimenti è quindi \(\Theta(m)\), e il
costo \emph{ammortizzato} di ciascuno è \(\Theta(1)\), anche se
una singola operazione che provoca il ridimensionamento costa
\(\Theta(n)\).

Inserire un elemento all'indice \(i\) richiede di spostare verso
destra i \(n-i\) elementi successivi; rimuovere quello di indice
\(i\) richiede di spostare a sinistra i \(n-i-1\) elementi seguenti.
A capacità sufficiente, i costi sono rispettivamente
\(\Theta(n-i+1)\) e \(\Theta(n-i)\). L'inserimento in coda è
un caso particolare, al quale si applica il costo ammortizzato
precedente.

\section{Liste concatenate}
\label{sec:fi-liste-concatenate}

Una \emph{lista concatenata} (\emph{linked list}) è composta da nodi,
ciascuno contenente un elemento e uno o più riferimenti ad altri nodi.
Nella \emph{lista semplicemente concatenata} ogni nodo contiene un
riferimento \texttt{next} al successivo; un riferimento \texttt{head}
individua il primo nodo e l'ultimo ha \texttt{next} nullo. I nodi
non devono occupare locazioni contigue.

\begin{figure}[htbp]
  \centering
  \begin{tikzpicture}[unipd/diagram,x=1cm,y=1cm]
  \node[anchor=east,font=\small] at (-4.4,0.8) {Array};
  \foreach \x/\v in {-3.6/A,-2.15/B,-0.7/C,0.75/D} {
    \node[unipd/node,minimum width=1cm,minimum height=.65cm] at (\x,0.8) {\v};
  }
  \node[anchor=east,font=\small] at (-4.4,-0.9) {Lista concatenata};
  \node[font=\scriptsize] (head) at (-4.1,-1.6) {\texttt{head}};
  \foreach \x/\v/\id in {-3.1/A/a,-1.2/B/b,0.7/C/c} {
    \node[unipd/node,minimum width=1.35cm,minimum height=.65cm] (\id) at (\x,-0.9) {\v\,\(\mid\)\,\scriptsize next};
  }
  \draw[unipd/directed edge] (head) -- (a.south);
  \draw[unipd/directed edge] (a) -- (b);
  \draw[unipd/directed edge] (b) -- (c);
  \node[font=\scriptsize] (null) at (2.35,-0.9) {\texttt{None}};
  \draw[unipd/directed edge] (c) -- (null);
\end{tikzpicture}

  \caption{Organizzazione contigua di un array e collegamenti tra i nodi di una lista semplicemente concatenata.}
  \label{fig:fi-array-lista-concatenata}
\end{figure}

Nella \cref{fig:fi-array-lista-concatenata} l'array consente di
raggiungere direttamente un indice, mentre nella lista concatenata
occorre seguire i riferimenti a partire da \texttt{head}. Per
raggiungere il nodo di indice \(i\) sono necessari \(i\) passaggi:
considerando anche l'accesso al nodo iniziale, il costo è
\(\Theta(i+1)\), e nel caso peggiore \(\Theta(n)\).

\subsection{Inserimento, rimozione e varianti}
\label{sec:fi-liste-concatenate-operazioni}

Inserire un nodo all'inizio richiede di collegarlo alla vecchia testa
e aggiornare \texttt{head}. Per rimuovere la testa è sufficiente
spostare \texttt{head} al nodo successivo. Entrambe le operazioni
costano \(\Theta(1)\), tenendo conto anche del caso di lista vuota.

Se è noto il riferimento a un nodo \(p\), un nuovo nodo \(q\) può
essere inserito \emph{dopo} \(p\) mediante i due aggiornamenti
\[
  q.\texttt{next}\gets p.\texttt{next},
  \qquad p.\texttt{next}\gets q.
\]
Analogamente, per rimuovere il nodo \emph{successivo} a \(p\), se
esiste, basta porre
\[
  p.\texttt{next}\gets p.\texttt{next.next}.
\]
L'aggiornamento dei collegamenti costa \(\Theta(1)\); se il nodo
di riferimento deve prima essere cercato, si aggiunge il costo
dell'attraversamento. In particolare, conoscere il nodo da
rimuovere non è in generale sufficiente per eliminarlo in tempo
costante da una lista semplicemente concatenata: può servire il
riferimento al suo predecessore. Un riferimento diretto
\texttt{tail} rende costante l'inserimento in fondo, ma non, da solo,
la rimozione dell'ultimo nodo.

In una \emph{lista doppiamente concatenata} ciascun nodo conserva
anche il riferimento \texttt{prev} al precedente, permettendo
l'attraversamento in entrambe le direzioni. In una \emph{lista
circolare}, invece, il collegamento dell'ultimo nodo riporta al
primo. Queste varianti modificano le operazioni disponibili e la
memoria necessaria per i riferimenti.

\begin{table}[htbp]
  \centering
  \caption{Costi delle principali operazioni, con \(n\) elementi e indici da zero.}
  \label{tab:fi-array-concatenata-costi}
  \small
  \begin{tabular}{lll}
    \hline
    Operazione & Array dinamico & Lista concatenata\\
    \hline
    Accesso all'indice \(i\) & \(\Theta(1)\) & \(\Theta(i+1)\)\\
    Inserimento in testa & \(\Theta(n+1)\) & \(\Theta(1)\)\\
    Rimozione in testa & \(\Theta(n)\) & \(\Theta(1)\)\\
    Inserimento in posizione nota & \(O(n)\) & \(\Theta(1)\)\\
    Inserimento in fondo & \(O(1)\) ammort. & \(\Theta(1)\) con \texttt{tail}\\
    \hline
  \end{tabular}
\end{table}

Nella riga relativa a una posizione nota, l'array riceve l'indice
di inserimento, mentre per la lista concatenata è già disponibile
il riferimento al nodo precedente.
Le liste concatenate sono adatte a inserimenti e rimozioni presso
nodi già individuati; gli array sono preferibili quando prevalgono
gli accessi per indice. La scelta deve considerare anche il costo
dei riferimenti e la località di memoria.

\section{Pila: politica LIFO}
\label{sec:fi-adt-pila}

\begin{definition}[Pila]
\label{def:fi-adt-pila}
Una \emph{pila} (\emph{stack}) è un ADT in cui inserimento,
rimozione e ispezione interessano soltanto la cima, secondo la
politica \emph{LIFO} (\emph{Last In, First Out}).
\end{definition}

Le operazioni sono \texttt{push(e)}, che inserisce \(e\),
\texttt{pop()}, che rimuove e restituisce la cima, e \texttt{top()},
che la restituisce senza rimuoverla. Le operazioni
\texttt{is\_empty()} e \texttt{size()} verificano rispettivamente
se la pila è vuota e quanti elementi contiene. \texttt{pop()} e
\texttt{top()} richiedono una pila non vuota: il comportamento in
caso contrario deve essere specificato dall'interfaccia.

\begin{example}[Ordine delle operazioni su una pila]
Le operazioni \texttt{push(10)} e \texttt{push(20)} collocano
\(20\) in cima. Una successiva \texttt{push(30)} porta in cima
\(30\). Tre chiamate a \texttt{pop()} restituiscono nell'ordine
\(30\), \(20\) e \(10\).
\end{example}

La pila delle chiamate della \cref{sec:fi-pila-ricorsione} segue
questa stessa politica. Altre applicazioni comprendono la gestione
delle operazioni di annullamento e il ritorno alle posizioni
precedenti durante una navigazione.

\subsection{Implementazioni della pila}
\label{sec:fi-adt-pila-implementazioni}

Con un array dinamico, la cima coincide con l'ultimo elemento
presente. \texttt{push} aggiunge in fondo e \texttt{pop} rimuove
l'ultimo elemento: entrambe le operazioni richiedono
\(O(1)\) tempo ammortizzato. \texttt{top}, \texttt{is\_empty}
e \texttt{size} richiedono \(\Theta(1)\).

Una realizzazione in Python utilizza direttamente queste operazioni
sulla lista:

\begin{unipdcode}[language=Python]
class Pila:
    def __init__(self):
        self._dati = []

    def push(self, elemento):
        self._dati.append(elemento)

    def pop(self):
        if self.is_empty():
            raise IndexError("Pila vuota")
        return self._dati.pop()

    def top(self):
        if self.is_empty():
            raise IndexError("Pila vuota")
        return self._dati[-1]

    def is_empty(self):
        return len(self._dati) == 0

    def size(self):
        return len(self._dati)
\end{unipdcode}


Con una lista semplicemente concatenata, la cima coincide con
\texttt{head}: inserimento, rimozione e lettura in testa richiedono
\(\Theta(1)\). Se si mantiene un contatore aggiornato a ogni modifica,
anche \texttt{size()} costa \(\Theta(1)\); senza contatore, richiede
\(\Theta(n)\) per attraversare la lista.

\section{Coda: politica FIFO}
\label{sec:fi-adt-coda}

\begin{definition}[Coda]
\label{def:fi-adt-coda}
Una \emph{coda} (\emph{queue}) è un ADT in cui gli inserimenti
avvengono in fondo e le rimozioni all'inizio, secondo la politica
\emph{FIFO} (\emph{First In, First Out}).
\end{definition}

Le operazioni principali sono \texttt{enqueue(e)}, che aggiunge un
elemento nel retro, \texttt{dequeue()}, che rimuove e restituisce
il primo, e \texttt{first()}, che lo ispeziona senza rimuoverlo.
Completano l'interfaccia \texttt{is\_empty()} e \texttt{size()}.
Rimozione e ispezione richiedono una coda non vuota. Le code di stampa
e l'elaborazione dei pacchetti di rete sono applicazioni della politica
FIFO.

\begin{example}[Ordine delle operazioni su una coda]
Dopo \texttt{enqueue(A)}, \texttt{enqueue(B)} e
\texttt{enqueue(C)}, la prima chiamata a \texttt{dequeue()}
restituisce \(A\), la seconda \(B\). L'elemento \(C\) rimane
all'inizio della coda.
\end{example}

\subsection{Implementazione mediante array circolare}
\label{sec:fi-coda-circolare}

Un array dinamico ordinario con rimozione dal primo indice
richiederebbe di spostare tutti gli elementi successivi, con costo
\(\Theta(n)\). Per evitarlo si usa un \emph{array circolare} con
capacità \(C>0\), indice \(f\) del primo elemento e numero
\(n\) di elementi presenti. La posizione logica \(i\) occupa
l'indice fisico
\[
  (f+i)\bmod C,
  \qquad 0\le i<n.
\]
In particolare, l'indice del prossimo inserimento è
\[
  b=(f+n)\bmod C.
\]
L'inserimento scrive in \(b\) e incrementa \(n\); la rimozione legge
l'elemento in \(f\), pone \(f\gets(f+1)\bmod C\) e decrementa \(n\).
Il resto permette di riutilizzare le posizioni liberate senza
spostare gli elementi, come illustra la
\cref{fig:fi-coda-circolare}.

\begin{figure}[htbp]
  \centering
  \begin{tikzpicture}[unipd/diagram,x=1cm,y=1cm]
  \foreach \i/\v in {0/E,1/F,2/{},3/{},4/C,5/D} {
    \node[unipd/node,minimum width=.85cm,minimum height=.68cm] (c\i) at ({(\i-2.5)*1.15},0) {\v};
    \node[font=\scriptsize] at ({(\i-2.5)*1.15},-0.65) {\i};
  }
  \node[font=\small] (front) at (1.725,1.25) {\texttt{front}=4};
  \draw[unipd/directed edge] (front) -- (c4.north);
  \node[font=\small] (back) at (-.575,1.25) {\texttt{back}=2};
  \draw[unipd/directed edge] (back) -- (c2.north);
  \node[font=\small] at (0,-1.22) {Ordine logico: \(C,D,E,F\)};
\end{tikzpicture}

  \caption{Coda circolare di capacità sei: \(f=4\), \(n=4\) e posizione del prossimo inserimento \(b=2\).}
  \label{fig:fi-coda-circolare}
\end{figure}

Se la capacità si esaurisce, il ridimensionamento ricopia gli
elementi nell'ordine logico a partire dal primo e aggiorna gli
indici. Con la strategia della
\cref{sec:fi-array-dinamici}, \texttt{enqueue} ha costo
\(O(1)\) ammortizzato, mentre \texttt{dequeue},
\texttt{first}, \texttt{is\_empty} e \texttt{size} hanno costo
\(\Theta(1)\).

La seguente implementazione in Python mantiene l'indice del primo
elemento e il numero di elementi presenti. Quando la capacità è
esaurita, raddoppia l'array ricopiando gli elementi nell'ordine logico.

\begin{unipdcode}[language=Python]
class CodaCircolare:
    def __init__(self):
        self._dati = [None] * 4
        self._inizio = 0
        self._dimensione = 0

    def is_empty(self):
        return self._dimensione == 0

    def size(self):
        return self._dimensione

    def enqueue(self, elemento):
        capacita = len(self._dati)
        if self._dimensione == capacita:
            nuovo = [None] * (2 * capacita)
            for i in range(self._dimensione):
                nuovo[i] = self._dati[(self._inizio + i) % capacita]
            self._dati = nuovo
            self._inizio = 0
            capacita *= 2
        indice = (self._inizio + self._dimensione) % capacita
        self._dati[indice] = elemento
        self._dimensione += 1

    def first(self):
        if self.is_empty():
            raise IndexError("Coda vuota")
        return self._dati[self._inizio]

    def dequeue(self):
        elemento = self.first()
        self._dati[self._inizio] = None
        self._inizio = (self._inizio + 1) % len(self._dati)
        self._dimensione -= 1
        return elemento
\end{unipdcode}


\subsection{Implementazione mediante lista concatenata}
\label{sec:fi-coda-concatenata}

Una lista semplicemente concatenata può mantenere sia il riferimento
\texttt{head}, corrispondente all'inizio, sia \texttt{tail},
corrispondente all'ultimo nodo. \texttt{enqueue} collega il nuovo
nodo dopo \texttt{tail}, mentre \texttt{dequeue} rimuove la testa.
Entrambe le operazioni hanno costo \(\Theta(1)\), purché si
aggiornino correttamente \texttt{head} e \texttt{tail} sia inserendo
il primo elemento sia rimuovendo l'ultimo.
Mantenendo un contatore, tutte le operazioni della coda hanno
costo \(\Theta(1)\). Senza riferimento \texttt{tail}, invece,
l'inserimento in fondo richiede \(\Theta(n)\).

\section{Coda doppia}
\label{sec:fi-adt-deque}

Una \emph{coda doppia} (\emph{deque}, da \emph{double-ended queue})
consente inserimenti, rimozioni e ispezioni a entrambe le estremità,
senza accesso arbitrario agli elementi interni. Le operazioni
\texttt{add\_first(e)} e \texttt{add\_last(e)} inseriscono rispettivamente
all'inizio e alla fine; \texttt{delete\_first()} e
\texttt{delete\_last()} rimuovono e restituiscono gli elementi
alle stesse estremità. \texttt{first()}, \texttt{last()},
\texttt{is\_empty()} e \texttt{size()} completano l'interfaccia.

La deque generalizza pila e coda: limitando l'accesso a una sola
estremità si ottiene la politica LIFO, mentre inserendo in fondo e
rimuovendo in testa si ottiene FIFO. Può essere realizzata con
un array circolare ridimensionabile oppure con una lista
doppiamente concatenata e riferimenti a entrambi gli estremi,
ottenendo inserimenti e rimozioni in \(O(1)\) ammortizzato
per l'array e \(\Theta(1)\) per la lista concatenata.
Fra le applicazioni figurano la gestione flessibile delle attività
e il confronto dei caratteri alle due estremità di una sequenza
per verificare se è palindroma.

\section{Lista come tipo di dato astratto}
\label{sec:fi-adt-lista}

L'ADT \emph{lista} descrive una sequenza in cui la posizione di
ciascun elemento è significativa. Gli elementi possono essere
letti, aggiornati, inseriti e rimossi a indici specificati. Le operazioni
\texttt{get(i)}, \texttt{set(i,e)}, \texttt{add(i,e)},
\texttt{append(e)} e \texttt{remove(i)} svolgono rispettivamente
lettura, sostituzione, inserimento, aggiunta in fondo e rimozione;
\texttt{is\_empty()} e \texttt{size()} descrivono lo stato del
contenitore. Nell'inserimento gli elementi successivi avanzano di
una posizione; la rimozione restituisce l'elemento eliminato.
Gli indici validi per l'accesso appartengono a \([0,n-1]\),
mentre per l'inserimento è ammesso anche \(i=n\).

La lista Python della \cref{sec:fi-python-liste} realizza queste
funzionalità, ma con un'interfaccia diversa: \texttt{insert(i, e)}
inserisce all'indice \(i\), \texttt{pop(i)} rimuove per indice,
mentre \texttt{remove(e)} elimina la prima occorrenza di \(e\).
La scelta tra array dinamico e lista concatenata dipende dal carico
di operazioni, come mostra la
\cref{tab:fi-array-concatenata-costi}: l'array privilegia l'accesso
per indice, mentre una struttura concatenata permette aggiornamenti
di collegamenti in tempo costante dopo aver individuato il nodo.

\section{Mappa}
\label{sec:fi-adt-mappa}

L'ADT \emph{mappa} (\emph{map}) associa valori a chiavi distinte.
Ogni chiave identifica al più un valore, mentre chiavi diverse
possono avere valori uguali. L'operazione \texttt{put(k,v)}
inserisce una nuova associazione o aggiorna il valore di una
chiave esistente; \texttt{get(k)} ne recupera il valore e
\texttt{remove(k)} elimina l'associazione. \texttt{contains(k)}
verifica la presenza della chiave; \texttt{keys()},
\texttt{values()} e \texttt{items()} espongono rispettivamente
chiavi, valori e associazioni. \texttt{is\_empty()} e
\texttt{size()} forniscono informazioni sul contenitore.
Per \texttt{get(k)} e \texttt{remove(k)}, il comportamento quando
la chiave è assente deve essere precisato dall'interfaccia.

\begin{example}[Conteggio delle frequenze]
Per ogni parola \(w\) letta in un testo, se \(w\) non è presente
nella mappa si registra l'associazione \(w\mapsto1\); altrimenti
si aggiorna il valore a \(\texttt{get}(w)+1\). Il risultato
associa a ciascuna parola il numero delle sue occorrenze.
\end{example}

La ricerca per chiave non implica che gli elementi siano ordinati
per chiave. Le complessità di accesso e aggiornamento dipendono
dall'implementazione della mappa: l'interfaccia astratta non
ne prescrive una specifica. In Python il tipo \texttt{dict}
fornisce un'implementazione di questo comportamento.

\section{Insieme}
\label{sec:fi-adt-insieme}

L'ADT \emph{insieme} (\emph{set}) contiene elementi distinti,
senza associare loro valori né richiedere un ordinamento.
\texttt{add(e)} inserisce \(e\) se assente,
\texttt{remove(e)} lo elimina e \texttt{contains(e)}
ne verifica l'appartenenza. L'interfaccia stabilisce il comportamento
di \texttt{remove(e)} quando l'elemento è assente.
Si considerano inoltre le operazioni
di unione, intersezione e differenza:
\[
  A\cup B,\qquad A\cap B,\qquad A\setminus B,
\]
nonché la verifica di inclusione \(A\subseteq B\), la consultazione
degli elementi e le operazioni \texttt{is\_empty()} e
\texttt{size()}. L'ordine con cui si elencano gli elementi non
fa parte della specifica astratta.

\begin{example}[Eliminazione dei duplicati]
Inserendo successivamente i valori \(4,2,2,8,6,2\) in un insieme,
si ottiene \(\{2,4,6,8\}\). Una successiva inserzione di \(2\)
non modifica il contenuto. Su insiemi di documenti, unione,
intersezione e differenza permettono di combinare risultati di
ricerca secondo la presenza di determinati termini.
\end{example}

L'ADT insieme è distinto dalla mappa: nel primo gli elementi
sono unici, nella seconda lo sono le chiavi e non necessariamente
i valori. Python offre il tipo \texttt{set}; anche in questo caso
le prestazioni delle singole operazioni sono una proprietà della
struttura dati adottata, non della definizione dell'ADT.
